\section{Fundamentação Teórica}

\begin{frame}{Modelo do membro superior}
  \begin{columns}[c]
    \begin{column}{0.40\textwidth}
      \centering
      \resizebox{!}{0.68\textheight}{\upperlimbschematic}

      {\scriptsize Adaptado de \textcite{Albuquerque:2025}.}
    \end{column}
    \begin{column}{0.58\textwidth}
      \begin{itemize}
        \item Modelo planar com 3 GdL: ombro ($\theta_1$), cotovelo ($\theta_2$) e punho ($\theta_3$) \cite{Gebai:2018}.
        \item Músculos e tendões como molas e amortecedores equivalentes; $k_3$ representa o acoplamento pelo bíceps.
      \end{itemize}
      \begin{equation*}
        \mathbf{J}\ddot{\bm{\theta}} + \mathbf{C}\dot{\bm{\theta}} + \mathbf{K}\bm{\theta} = \bm{\tau},
        \,\,
        \mathbf{K} =
        \begin{bmatrix}
          k_1 + k_3 & k_3       & 0   \\
          k_3       & k_2 + k_3 & 0   \\
          0         & 0         & k_4
        \end{bmatrix}
      \end{equation*}
      \begin{itemize}
        \item $\mathbf{C}$ tem a estrutura de $\mathbf{K}$, com $c_i \propto k_i$; $\mathbf{J}$ é cheia e acopla as três articulações.
        \item Linearização adequada se os parâmetros variam pouco no tempo \cite{Gebai:2020}.
      \end{itemize}
    \end{column}
  \end{columns}
\end{frame}

\begin{frame}{Torques aplicados e representação em espaço de estados}
  \begin{columns}[T]
    \begin{column}{0.5\textwidth}
      \begin{block}{Torque em um paciente com DP \cite{Qi:2024}}
        \vspace{-0.3cm}
        \begin{equation*}
          \bm{\tau} = \underbrace{\bm{\tau}_i}_{\text{patológico}} + \underbrace{\bm{\tau}_v}_{\text{voluntário}}
        \end{equation*}
      \end{block}
      \begin{itemize}
        \item A potência do tremor concentra-se em torno de uma frequência central \cite{Deuschl:2022}; usualmente, até 3 harmônicos \cite{Zamanian:2019,Albuquerque:2025}.
        \item O acoplamento já torna a dinâmica do punho multimodal: um harmônico por articulação.
      \end{itemize}
    \end{column}
    \begin{column}{0.46\textwidth}
      \begin{block}{Tremor}
        \begin{equation*}
          \bm{\tau}_i =
          \begin{bmatrix}
            A_1\sin(2\pi f_1 t) \\
            A_2\sin(2\pi f_2 t) \\
            A_3\sin(2\pi f_3 t)
          \end{bmatrix}
        \end{equation*}
        \centering
        \small
        $A_i = 1$~N$\cdot$m;\quad $f_1\ f_2\ f_3 = 3{,}59\ 5{,}30\ 14{,}35$~Hz
      \end{block}
      \begin{block}{Espaço de estados, $\mathbf{x} = [\bm{\theta}^\top\ \dot{\bm{\theta}}^\top]^\top$}
        \vspace{-0.3cm}
        \begin{align*}
          \dot{\mathbf{x}} &= \mathbf{A}\mathbf{x} + \mathbf{B}\bm{\tau}, \qquad
          \mathbf{y} = \mathbf{C}_\text{ss}\mathbf{x} \\[1mm]
          \mathbf{A} &=
          \begin{bmatrix}
            \mathbf{0} & \mathbf{I} \\
            -\mathbf{J}^{-1}\mathbf{K} & -\mathbf{J}^{-1}\mathbf{C}
          \end{bmatrix},
          \quad
          \mathbf{B} =
          \begin{bmatrix}
            \mathbf{0} \\
            \mathbf{J}^{-1}
          \end{bmatrix} \\[1mm]
          \mathbf{C}_\text{ss} &=
          \begin{bmatrix}
            \mathbf{I} & \mathbf{0}
          \end{bmatrix}
        \end{align*}
      \end{block}
    \end{column}
  \end{columns}
\end{frame}

\begin{frame}{Controle por rejeição ativa de perturbações baseado em erro}
  \begin{columns}[T]
    \begin{column}{0.48\textwidth}
      \begin{itemize}
        \item ADRC: estima e rejeita, em tempo real, a \alert{perturbação total}, isto é, incertezas internas e perturbações externas \cite{Guo:2016}:
          \begin{equation*}
            y^{(n)} = f(y, \dot{y}, \dots, y^{(n-1)}, t) + b_0 u.
          \end{equation*}
        \item EADRC: mesma estrutura no erro de rastreamento $e = r - y$, de implementação mais simples \cite{Madonski:2019}:
          \begin{equation*}
            e^{(n)} = f_e - b_0 u,
            \qquad
            f_e = r^{(n)} - f(\cdot).
          \end{equation*}
      \end{itemize}
    \end{column}
    \begin{column}{0.48\textwidth}
      ESO, com $z_i = e^{(i-1)}$ ($i \leq n$) e $z_{n+1} = f_e$:
      \begin{equation*}
        \begin{cases}
          \dot{\hat{z}}_i = \hat{z}_{i+1} + \beta_i (e - \hat{z}_1), \quad i < n \\
          \dot{\hat{z}}_n = \hat{z}_{n+1} + \beta_n (e - \hat{z}_1) - b_0 u \\
          \dot{\hat{z}}_{n+1} = \beta_{n+1} (e - \hat{z}_1)
        \end{cases}
      \end{equation*}
      Lei de controle e dinâmica resultante do erro:
      \begin{align*}
        u &= \tfrac{1}{b_0}\big(\textstyle\sum_{i=1}^{n} \kappa_i \hat{z}_i + \hat{z}_{n+1}\big) \\
        \implies\ e^{(n)} &\approx -\textstyle\sum_{i=1}^{n} \kappa_i e^{(i-1)}
      \end{align*}
    \end{column}
  \end{columns}

  \begin{alertblock}{Cuidado}
    \small
    Usualmente, limita-se a magnitude ou a taxa do sinal de controle 
    para evitar instabilidade do observador\cite{Herbst:2025}.
  \end{alertblock}
\end{frame}
